\section*{II Une atmosphère électrique}

L'atmosphère est assimilée à un plasma dans lequel le courant électrique est essentiellement dû à des déplacements d'électrons non relativistes. La densité d'électrons mobiles par unité de volume est notée $n$, la masse d'un électron $m_e$, sa charge $e$.

On étudie, dans un premier temps, la propagation d'une onde électromagnétique plane progressive monochromatique de pulsation $\omega$, de vecteur d'onde $\vec{k} = k\vec{e}_z$, polarisée rectilignement et décrite par son vecteur champ électrique à l'instant $t$, $\vec{E} = \vec{E}_0 \exp i(\omega t - kz)$ avec $\vec{E}_0 \perp \vec{e}_z$ et $\|\vec{e}_z\| = 1$.

L'atmosphère est caractérisée par la permittivité diélectrique du vide $\varepsilon_0$ et par la perméabilité magnétique du vide $\mu_0$. Dans cette étude, on néglige les chocs que subissent les électrons et leur poids. La seule force agissant sur les électrons est la force de LORENTZ.

\begin{itemize}
    \item[13.] Montrer que la contribution magnétique de la force de LORENTZ sur l'électron est négligeable devant la contribution électrique de cette même force.
    \item[14.] Justifier le fait que la contribution des ions au courant électrique est négligeable devant celle des électrons. Déterminer la relation qui existe entre la représentation complexe du champ électrique $\vec{E}$ de l'onde électromagnétique et celle de la densité volumique de courant $\vec{j}$ qui résulte de l'interaction entre l'électron et $\vec{E}$. Donner l'expression de la conductivité complexe du plasma $\gamma_p$ et commenter cette expression en terme de puissance transférée.
    \item[15.] Établir l'équation de propagation du champ électrique $\vec{E}$. En déduire la relation de dispersion vérifiée par le vecteur d'onde $k$.
    \item[16.] Pour quelle fréquence limite $f_p$, une onde envoyée depuis la surface de la Terre ne peut plus se propager lorsqu'elle rencontre un milieu où la densité volumique d'électrons est $n$ ? On détaillera le raisonnement permettant d'obtenir cette limite. Effectuer l'application numérique pour une densité $n = 10^{11}\text{ m}^{-3}$ caractéristique des sylphes.
    \item[17.] L'image de la figure 6 montre la localisation et le nombre d'échos radar qui reviennent au niveau de l'émetteur. Analyser et commenter cette carte.
\end{itemize}

\TLEFigure{6}{Carte d'échos radar à 2,2 MHz depuis le sol pour l'évènement enregistré à 13h39 le 15 octobre 1994 dans la station de Khorogo (Côte d'Ivoire) -- figure tirée de l'article \textit{HF echoes from ionization potentially produced by high-altitude discharges}, R.A. Roussel-Dupré \& E. Blanc, Journal of geophysical research, Vol. 102, n$^\circ$A3, p. 4613-4622, 1997}
